% Generated draft: requires mathematical and source review.
\documentclass[11pt]{article}
\usepackage[margin=25mm]{geometry}
\usepackage{fontspec}
\setmainfont{Latin Modern Roman}
\usepackage{amsmath,amssymb,mathtools,unicode-math}
\setmathfont{Latin Modern Math}
\usepackage{enumitem,needspace,xurl,hyperref}
\setlength{\parindent}{0pt}
\setlength{\parskip}{5pt}
\setlength{\emergencystretch}{3em}
\newcommand{\N}{\mathbb{N}}
\newcommand{\Z}{\mathbb{Z}}
\newcommand{\Q}{\mathbb{Q}}
\newcommand{\R}{\mathbb{R}}
\newcommand{\C}{\mathbb{C}}
\newcommand{\F}{\mathbb{F}}
\newcommand{\A}{\mathbb{A}}
\newcommand{\PP}{\mathbb{P}}
\newcommand{\E}{\mathbb{E}}
\newcommand{\eps}{\varepsilon}
\newcommand{\dd}{\,\mathrm{d}}
\begin{document}
\section*{UnsolvedMath ID: 30004655}
\textbf{Source problem code:} catalogue entry
\subsection{Definitions and mathematical statement}
Let \(d\) be the spatial dimension, with no restriction on positive integer dimensions specified in the supplied record. Let \(a:\R^d\to\R^{d\times d}\) be a uniformly elliptic matrix field that is \(\Z^d\)-periodic, meaning that
\[
a(y+k)=a(y)\qquad (k\in\Z^d).
\]
The record does not specify the regularity, symmetry, or precise boundedness convention for \(a\). Let \(W:\R\to\R\) be the stated double-well potential, with global minima at \(-1\) and \(1\); in particular,
\[
W(-1)=W(1)=\inf_{s\in\R}W(s).
\]
Here \(W'\) denotes the derivative appearing in the equation; the source record does not give further regularity, growth, or nondegeneracy assumptions on \(W\).

For each \(\varepsilon>0\), let \(u_\varepsilon:\R^d\times[0,\infty)\to\R\) solve the initial-value problem
\[
\begin{cases}
\partial_tu_\varepsilon-\operatorname{div}_x\!\left(a(x/\varepsilon)D_xu_\varepsilon\right)+\varepsilon^{-2}W'(u_\varepsilon)=0,
& (x,t)\in\R^d\times(0,\infty),\\
u_\varepsilon(x,0)=u_0(x),&x\in\R^d.
\end{cases}
\]
Here \(D_x\) and \(\operatorname{div}_x\) denote the spatial gradient and divergence. The datum \(u_0\) is written without dependence on \(\varepsilon\), as in the supplied statement. Its admissible class and the meaning of solution and initial trace are unspecified. The spatial domain is the whole space; no boundary condition or behavior at spatial infinity is supplied.

\textbf{Problem.} Determine the limiting behavior of \(u_\varepsilon\) as \(\varepsilon\to0^+\). The record does not specify a convergence topology, a class of admissible initial data, whether convergence is sought for the entire family or only subsequences, or a proposed limiting evolution. These omissions prevent a fully quantified convergence assertion. The title indicates a sharp-interface target, but no limiting interface or interface solution concept is defined in the supplied statement.

The original statement additionally asks for the limit for a periodic mean-curvature-type equation. That equation and its hypotheses are absent from the supplied record, so this additional part cannot be formulated precisely here.
\subsection{Short English statement}
Determine the small-\(\varepsilon\) behavior of the stated phase-transition equation in a periodic uniformly elliptic medium, with fixed initial datum \(u_0\). The source leaves the admissible data, solution concept, convergence mode, and limiting interface law unspecified. It also mentions a second, unspecified periodic mean-curvature-type equation.
\subsection{Sources}
\begin{itemize}
\item Local UnsolvedMath registry: \texttt{data/unsolvedmath/problems.json}, with \texttt{catalog/problems.json} as the supplementary registry; supplied permanent ID \(30004655\), exact source problem code \texttt{catalogue entry}, title \emph{Sharp-Interface Limits of Periodic Phase-Transition Equations}. The supplied collection is \emph{Oberwolfach Reports Open Problems}. No individual proposer is identified; the supplied proposed year is 2021.
\item Supplied report citation: \emph{Homogenization Theory: Periodic and Beyond} (2022), Oberwolfach Reports, pp.~791--852, DOI \texttt{10.4171/owr/2021/14}: \url{https://doi.org/10.4171/owr/2021/14}. The same reference is listed among the literature sources as \emph{Oberwolfach Report 14/2021}. The record describes Oberwolfach Reports as documenting mathematical workshops and their open-problem discussions. The report contents have not been independently consulted for this draft.
\item Supplied literature source: Morfe, \emph{Homogenization of Interface Motions}, \url{https://doi.org/10.6082/uchicago.4837}.
\item Supplied literature source: Morfe, \emph{Allen--Cahn with periodic mobility}, \url{https://arxiv.org/abs/2011.15012}.
\item Supplied literature assessment, checked on \texttt{2026-08-22}: the record labels the problem \texttt{partially\_solved}. It reports that homogenized interface laws and sharp-interface limits are proved for laminar media, restricted initial data, and periodic mobility/dissipation models, while the full periodic uniformly elliptic phase-transition problem as stated remains unresolved in all dimensions and data. This is supplied dated curation derived from the OWR material, not an independently verified literature assessment; its time-sensitive status and cited claims require independent verification.
\item The supplied statement status is \texttt{corrected\_verified}, accompanied by a claim that the clean statement was checked against the official report and preserves its hypotheses, notation, and target. This is a registry verification claim, not verification performed for this draft; the missing mean-curvature-type equation remains unresolved.
\end{itemize}
\end{document}
